\hdB{配点一覧（100点満点／200点換算併記）}
\begin{center}\small
\begin{tabular}{@{}clcr@{}}
\toprule
\textbf{大問} & \textbf{分野・主題} & \textbf{小問の配点} & \textbf{計} \\
\midrule
1 & 数学I・A・II 小問集合（対称式・円内接四角形・条件付き確率・対数と桁数） & $5+5+5+5$ & 20 \\
2 & 数学II 微分法・積分法（3次関数の対称性・接線・面積比・線分長の最大） & $4+5+5+6$ & 20 \\
3 & 数学II 三角関数・高次方程式（3倍角・3次方程式の解の配置・解の個数分類） & $4+5+5+6$ & 20 \\
4 & 数学B 数列（$S_{n}$ と $a_{n}$ の漸化式・指数型漸化式・$(等差)\times(等比)$の和） & $4+5+5+6$ & 20 \\
5 & 数学C 空間ベクトル（等面四面体の垂線と体積・対辺中点の直交・外接球の半径） & $4+5+5+6$ & 20 \\
\bottomrule
\end{tabular}
\end{center}

\hdB{自己採点 成績評価基準（志望校判定・偏差値目安・到達水準）}
\par\vspace{1.0mm}\noindent
{\footnotesize 本模試の自己採点得点（100点満点／2倍で200点換算）に対応する実力到達度・合格目安・今後の学習課題です。}
\par\vspace{1.5mm}
\begin{center}\footnotesize
\renewcommand{\arraystretch}{1.28}
\begin{tabular}{@{}cclp{26em}@{}}
\toprule
\textbf{得点（100点）} & \textbf{換算点（200点）} & \textbf{判定・偏差値目安} & \textbf{実力評価と今後の学習課題} \\
\midrule
$85 \sim 100$ 点 & $170 \sim 200$ 点 & \textbf{Sランク}（偏差値 68〜） & \textbf{東大・京大・難関国公立・早慶合格水準}。数I・II・B・Cの記述答案作成力は完成の域。計算ミス・論理の穴が極めて少なく，難関大実戦演習（第3セット）へ進む実力があります。 \\
$70 \sim 84$ 点 & $140 \sim 168$ 点 & \textbf{Aランク}（偏差値 60〜67） & \textbf{旧帝大・上位国公立大・上位私大理工合格水準}。典型大問を最後まで書き切る記述力があります。減点項目（真数条件・等号成立条件・垂直条件）を洗い出して盤石にしましょう。 \\
$50 \sim 69$ 点 & $100 \sim 138$ 点 & \textbf{Bランク}（偏差値 52〜59） & \textbf{地方国公立大・中堅私大記述合格水準}。大問1の小問集合や各大問の(1)(2)は確保できています。後半の小問で方針点を確実に拾う記述戦略を磨きましょう。 \\
$35 \sim 49$ 点 & $70 \sim 98$ 点 & \textbf{Cランク}（偏差値 45〜51） & \textbf{基礎定着レベル}。基本公式・定義（3次関数の極値，空間ベクトルの内積，漸化式の特性方程式など）の証明と適用条件を教科書レベルから再確認が必要です。 \\
$35$ 点未満 & $70$ 点未満 & \textbf{Dランク}（偏差値 〜44） & \textbf{要・基礎徹底復習}。各大問の(1)の標準解法を網羅系問題集（傍用問題集・青チャート・Focus Gold等）で復習し，基礎公式を自力で再現できるようにしましょう。 \\
\bottomrule
\end{tabular}
\end{center}

\rubricintro{対数の真数条件，$\theta$ と $t=\cos\theta$ の $1$ 対 $1$ 対応，漸化式における $n=1$ の確認，空間ベクトルにおける垂直条件の独立な $2$ 直線への適用など，論理の根幹に関わる条件を欠いた答案は減点の対象とする。}

\hdA{大問 1（20点）小問集合}

\Rh{1}{1}{5}
\Seikai{$x+y=4$，$xy=1$，$x^{3}+y^{3}=52$，$x^{5}+y^{5}=724$}
\begin{saiten}
\Sp{分母を有理化して $x=2+\sqrt{3}$，$y=2-\sqrt{3}$ を求めた}{+1}
\Sp{基本対称式 $x+y=4$，$xy=1$ をともに正しく求めた}{+1}
\Sp{$x^{3}+y^{3}=(x+y)^{3}-3xy(x+y)=52$ を正しく求めた}{+1}
\Sp{$x^{2}+y^{2}=14$ を経て $x^{5}+y^{5}=(x^{2}+y^{2})(x^{3}+y^{3})-x^{2}y^{2}(x+y)=724$ を正しく求めた}{+2}
\Ssep
\Sm{$x^{5}+y^{5}=(x^{2}+y^{2})(x^{3}+y^{3})$ として余分な項 $x^{2}y^{2}(x+y)=4$ を引き忘れた}{$-1$}
\end{saiten}
\betsunote{別解（漸化式 $S_{n+2}=4S_{n+1}-S_{n}$ を用いて順に計算する解法）でも同配点：$x+y, xy$ に $+2$，$S_{3}=52$ に $+1$，$S_{5}=724$ に $+2$。}

\Rh{1}{2}{5}
\Seikai{$\cos\angle\Pt{ABC}=-\dfrac{3}{7}$，$S=6\sqrt{10}$，$\Pt{AC}=\dfrac{\sqrt{1729}}{7}$}
\begin{saiten}
\Sp{円に内接する四角形の対角の性質から $\cos\angle\Pt{ADC}=-\cos\angle\Pt{ABC}$ を用いて余弦定理を $2$ つ立てた}{+1}
\Sp{$\Pt{AC}^{2}$ を消去して $\cos\angle\Pt{ABC}=-\dfrac{3}{7}$ を求めた}{+1}
\Sp{$\sin\angle\Pt{ABC}=\dfrac{2\sqrt{10}}{7}$ を求め，$S=\triangle\Pt{ABC}+\triangle\Pt{ADC}=6\sqrt{10}$ を求めた}{+2}
\Sp{$\Pt{AC}^{2}=\dfrac{247}{7}$ から $\Pt{AC}=\dfrac{\sqrt{1729}}{7}$ を求めた}{+1}
\Ssep
\Sm{$\cos(\pi-\theta)=\cos\theta$ として符号を誤った}{$-3$}
\end{saiten}

\Rh{1}{3}{5}
\Seikai{$p=\dfrac{209}{864}$，$q=\dfrac{128}{209}$}
\begin{saiten}
\Sp{硬貨 $\Pt{A}$ を選び表が $2$ 回出る確率 $P(A\cap E)=\dfrac{1}{3}\cdot\Cb{3}{2}\left(\dfrac{2}{3}\right)^{2}\left(\dfrac{1}{3}\right)=\dfrac{4}{27}$ を求めた}{+1}
\Sp{硬貨 $\Pt{B}$ を選び表が $2$ 回出る確率 $P(B\cap E)=\dfrac{2}{3}\cdot\Cb{3}{2}\left(\dfrac{1}{4}\right)^{2}\left(\dfrac{3}{4}\right)=\dfrac{3}{32}$ を求めた}{+1}
\Sp{排反事象の和として $p=\dfrac{4}{27}+\dfrac{3}{32}=\dfrac{209}{864}$ を求めた}{+1}
\Sp{条件付き確率の定義 $q=\dfrac{P(A\cap E)}{P(E)}$ に代入して $q=\dfrac{128}{209}$ を求めた}{+2}
\Ssep
\Sm{反復試行の係数 $\Cb{3}{2}=3$ を掛け忘れた}{$-2$}
\end{saiten}

\Rh{1}{4}{5}
\Seikai{$1<x<4,\ 4<x<7$，$32$ 桁，最高位の数字は $1$}
\begin{saiten}
\Sp{真数条件 $1<x<7$ を明記した}{+1}
\Sp{$(x-1)(7-x)<9 \iff (x-4)^{2}>0$ を導き，$1<x<4,\ 4<x<7$ を求めた}{+2}
\Sp{$\log_{10}6^{40}=31.124$ から $32$ 桁であることを導いた}{+1}
\Sp{$\log_{10}1<0.124<\log_{10}2$ から最高位の数字が $1$ であることを示した}{+1}
\Ssep
\Sm{真数条件 $1<x<7$ の確認がない}{$-1$}
\Sm{$(x-4)^{2}>0$ から $x=4$ を除き忘れ，$1<x<7$ と答えた}{$-1$}
\end{saiten}

\hdA{大問 2（20点）3次関数の対称性・接線・面積比}

\Rh{2}{1}{4}
\Seikai{極大値 $16$（$x=-1$），極小値 $-16$（$x=3$），共有点 $(1-2\sqrt{3}, 0),\ (1, 0),\ (1+2\sqrt{3}, 0)$}
\begin{saiten}
\Sp{$f'(x)=3(x+1)(x-3)$ を求め，増減と極大値 $16$（$x=-1$）・極小値 $-16$（$x=3$）を求めた}{+2}
\Sp{$f(x)=(x-1)(x^{2}-2x-11)$ と因数分解し，$x$ 軸との共有点 $3$ 個をすべて求めた}{+1}
\Sp{極値と共有点を反映した曲線 $C$ の概形を正しく描いた}{+1}
\Ssep
\Sm{極値をとる $x$ の値を明記していない}{$-1$}
\end{saiten}

\Rh{2}{2}{5}
\Seikai{$f(1+h)+f(1-h)=0$ の証明，接線 $\ell: y=-12x+12$}
\begin{saiten}
\Sp{$f(1+h)=h^{3}-12h$ および $f(1-h)=-h^{3}+12h$ を計算した}{+2}
\Sp{$f(1+h)+f(1-h)=0$ を示し，点 $\Pt{A}(1, 0)$ に関する点対称性を結論づけた}{+1}
\Sp{$f'(1)=-12$ を求め，接線 $\ell: y=-12x+12$ を求めた}{+2}
\Ssep
\Sm{特定の $h$ の値（$h=1, 2$ など）の代入のみで一般の $h$ の証明になっていない}{$-2$}
\end{saiten}

\Rh{2}{3}{5}
\Seikai{$\Pt{Q}(5, 16)$，$S_{1}=108$}
\begin{saiten}
\Sp{接線 $m$ の方程式 $y=16$ を求め，$16-f(x)=(5-x)(x+1)^{2}$ から $\Pt{Q}(5, 16)$ を求めた}{+2}
\Sp{面積を $S_{1}=\dsp\int_{-1}^{5}\{16-f(x)\}\,dx$ と正しく立式した}{+1}
\Sp{定積分を計算して $S_{1}=108$ を求めた}{+2}
\Ssep
\Sm{曲線と直線の上下関係の確認なく被積分関数の符号を逆にした}{$-1$}
\end{saiten}

\Rh{2}{4}{6}
\Seikai{$t=3$ のとき最大値 $32$，$S_{2}:S_{1}=16:27$}
\begin{saiten}
\Sp{線分 $\Pt{MN}$ の長さを $L(t)=16-f(t)$ と表した}{+1}
\Sp{$-1<t<5$ において $f(t)$ が $t=3$ で極小値 $-16$ をとることから，$t=3$ で最大値 $32$ をとることを示した}{+2}
\Sp{$S_{2}=\dsp\int_{-1}^{3}\{16-f(x)\}\,dx=64$ を正しく求めた}{+2}
\Sp{比 $S_{2}:S_{1}=64:108=16:27$ を求めた}{+1}
\Ssep
\Sm{比を最も簡単な整数の比に約分していない}{$-1$}
\end{saiten}

\hdA{大問 3（20点）三角関数・高次方程式の解の配置}

\Rh{3}{1}{4}
\Seikai{$\cos 3\theta=4t^{3}-3t$，$\theta=\dfrac{\pi}{9},\ \dfrac{5\pi}{9},\ \dfrac{7\pi}{9}$}
\begin{saiten}
\Sp{加法定理と $2$ 倍角の公式を用いて $\cos 3\theta=4t^{3}-3t$ を導出した}{+2}
\Sp{$0<3\theta<3\pi$ の範囲から $\cos 3\theta=\dfrac{1}{2}$ を満たす $\theta=\dfrac{\pi}{9},\ \dfrac{5\pi}{9},\ \dfrac{7\pi}{9}$ をすべて求めた}{+2}
\Ssep
\Sm{公式の結果のみを記し，加法定理からの導出過程がない}{$-1$}
\Sm{$\theta=\dfrac{7\pi}{9}$ を落とした}{$-1$}
\end{saiten}

\Rh{3}{2}{5}
\Seikai{$\alpha=\cos\dfrac{7\pi}{9},\ \beta=\cos\dfrac{5\pi}{9},\ \gamma=\cos\dfrac{\pi}{9}$，基本対称式 $0,\ -\dfrac{3}{4},\ \dfrac{1}{8}$，$\alpha^{2}+\beta^{2}+\gamma^{2}=\dfrac{3}{2}$}
\begin{saiten}
\Sp{$8t^{3}-6t-1=0 \iff 4t^{3}-3t=\dfrac{1}{2}$ と変形し，$3$ 解を $\cos\theta$ で表した}{+1}
\Sp{$0<\theta<\pi$ での $\cos\theta$ の単調減少性から大小を比べ，$\alpha, \beta, \gamma$ を正しく対応させた}{+1}
\Sp{解と係数の関係から $\alpha+\beta+\gamma=0$，$\alpha\beta+\beta\gamma+\gamma\alpha=-\dfrac{3}{4}$，$\alpha\beta\gamma=\dfrac{1}{8}$ を求めた}{+2}
\Sp{$\alpha^{2}+\beta^{2}+\gamma^{2}=\dfrac{3}{2}$ を求めた}{+1}
\Ssep
\Sm{$\alpha, \beta, \gamma$ の大小順序を逆にした（$\cos\dfrac{\pi}{9}$ を最小とした）}{$-1$}
\end{saiten}

\Rh{3}{3}{5}
\Seikai{極大値 $1$（$t=-\dfrac{1}{2}$），極小値 $-3$（$t=\dfrac{1}{2}$），各不等式の証明}
\begin{saiten}
\Sp{$g'(t)=6(2t+1)(2t-1)$ から増減と極大値 $1$（$t=-\dfrac{1}{2}$）・極小値 $-3$（$t=\dfrac{1}{2}$）を求めた}{+2}
\Sp{$g\!\left(-\dfrac{\sqrt{3}}{2}\right)=-1<0$，$g(0)=-1<0$，$g\!\left(\dfrac{\sqrt{3}}{2}\right)=-1<0$，$g(1)=1>0$ を計算した}{+2}
\Sp{各区間における単調性と端点の符号（中間値の定理）から $3$ つの不等式を示した}{+1}
\Ssep
\Sm{端点の符号の確認なく結論のみを述べた}{$-2$}
\end{saiten}
\betsunote{別解（$0<\theta<\pi$ における $\cos\theta$ の単調減少性と角度の大小比較による証明）でも同配点：極値に $+2$，角度の比較と単調性による証明に $+3$。}

\Rh{3}{4}{6}
\Seikai{$-1<a\leqq\dfrac{1}{3}$ のとき $3$ 個，$a\leqq -1$ または $a>\dfrac{1}{3}$ のとき $1$ 個}
\begin{saiten}
\Sp{$0\leqq\theta\leqq\pi$ において $\theta$ と $t=\cos\theta$（$-1\leqq t\leqq 1$）が $1$ 対 $1$ に対応することを明記した}{+1}
\Sp{方程式を $t\{4t^{2}-3(a+1)\}=0$ と因数分解し，$t=0$ がつねに解であることを示した}{+1}
\Sp{$t^{2}=\dfrac{3(a+1)}{4}$ が $0<t^{2}\leqq 1$ を満たす条件 $-1<a\leqq\dfrac{1}{3}$ を導いた}{+2}
\Sp{境界 $a=-1$（$t=0$ で重解）および $a=\dfrac{1}{3}$（$t=\pm 1$ を含む）を正しく処理して個数を分類した}{+2}
\Ssep
\Sm{$a=-1$ のときに $t=0$ の重複を見落とし $3$ 個とした}{$-1$}
\Sm{$a=\dfrac{1}{3}$ の等号を含め忘れた}{$-1$}
\end{saiten}

\hdA{大問 4（20点）和 $S_{n}$ を含む漸化式・指数型漸化式}

\Rh{4}{1}{4}
\Seikai{$a_{1}=2$，$a_{2}=14$，$a_{3}=58$，$a_{n+1}=3a_{n}+2^{n+2}$}
\begin{saiten}
\Sp{$n=1$ で $S_{1}=a_{1}$ を用いて $a_{1}=2$ を求めた}{+1}
\Sp{$a_{n+1}=S_{n+1}-S_{n}$ から漸化式 $a_{n+1}=3a_{n}+2^{n+2}$ を導いた}{+2}
\Sp{$a_{2}=14$，$a_{3}=58$ をともに正しく求めた}{+1}
\Ssep
\Sm{$2^{n+2}-2^{n+1}$ の整理を誤った}{$-2$}
\end{saiten}

\Rh{4}{2}{5}
\Seikai{$b_{n+1}=\dfrac{3}{2}b_{n}+2$，$b_{n}=5\left(\dfrac{3}{2}\right)^{n-1}-4$，$a_{n}=10\cdot 3^{n-1}-2^{n+2}$}
\begin{saiten}
\Sp{両辺を $2^{n+1}$ で割り，$b_{n+1}=\dfrac{3}{2}b_{n}+2$ を導いた}{+2}
\Sp{$b_{1}=1$ と特性方程式から $b_{n}=5\left(\dfrac{3}{2}\right)^{n-1}-4$ を求めた}{+2}
\Sp{$a_{n}=2^{n}b_{n}=10\cdot 3^{n-1}-2^{n+2}$ を求めた}{+1}
\Ssep
\Sm{初項 $b_{1}=\dfrac{a_{1}}{2}=1$ を $b_{1}=2$ と誤った}{$-2$}
\end{saiten}

\Rh{4}{3}{5}
\Seikai{$S_{n}=5\cdot 3^{n}-2^{n+3}+3$，最小の自然数 $n=4$}
\begin{saiten}
\Sp{$S_{n}=\dfrac{3}{2}a_{n}-2^{n+1}+3$（または等比数列の和）から $S_{n}=5\cdot 3^{n}-2^{n+3}+3$ を求めた}{+2}
\Sp{$a_{n}>10\cdot 2^{n}$ を $\left(\dfrac{3}{2}\right)^{n-1}>\dfrac{14}{5}=2.8$ に帰着させた}{+2}
\Sp{$n=3$（$2.25$）と $n=4$（$3.375$）を比較して最小の $n=4$ を求めた}{+1}
\Ssep
\Sm{$n=3$ での不成立を示さずに $n=4$ のみを確認した}{$-1$}
\end{saiten}

\Rh{4}{4}{6}
\Seikai{$W_{n}=\dfrac{9}{4}-\dfrac{2n+3}{4\cdot 3^{n-1}}$}
\begin{saiten}
\Sp{$c_{k}=3^{k-1}$ を求め，$W_{n}=\dsp\sum_{k=1}^{n}\frac{k}{3^{k-1}}$ と表した}{+1}
\Sp{$W_{n}-\dfrac{1}{3}W_{n}$ を作ってずらして引く方針を示した}{+2}
\Sp{$\dfrac{2}{3}W_{n}=\dfrac{3}{2}\left(1-\dfrac{1}{3^{n}}\right)-\dfrac{n}{3^{n}}$ まで正しく整理した}{+2}
\Sp{$W_{n}=\dfrac{9}{4}-\dfrac{2n+3}{4\cdot 3^{n-1}}$ を導いた}{+1}
\Ssep
\Sm{ずらして引いた際の末項 $-\dfrac{n}{3^{n}}$ の符号または指数を誤った}{$-2$}
\end{saiten}

\hdA{大問 5（20点）空間ベクトルと等面四面体の外接球}

\Rh{5}{1}{4}
\Seikai{$\vec{a}\cdot\vec{b}=2$，$\vec{a}\cdot\vec{c}=2$，$\vec{b}\cdot\vec{c}=3$，$\triangle\Pt{OBC}=2$}
\begin{saiten}
\Sp{$|\vec{b}-\vec{a}|^{2}=5$ 等から $\vec{a}\cdot\vec{b}=2$，$\vec{a}\cdot\vec{c}=2$，$\vec{b}\cdot\vec{c}=3$ をすべて求めた}{+2}
\Sp{$\triangle\Pt{OBC}=\dfrac{1}{2}\sqrt{|\vec{b}|^{2}|\vec{c}|^{2}-(\vec{b}\cdot\vec{c})^{2}}=2$ を求めた}{+2}
\Ssep
\Sm{内積の展開における $-2\vec{a}\cdot\vec{b}$ の符号を誤った}{$-2$}
\end{saiten}

\Rh{5}{2}{5}
\Seikai{$\Vec{OH}=\dfrac{1}{4}\vec{b}+\dfrac{1}{4}\vec{c}$，$\Pt{AH}=\sqrt{3}$，$V=\dfrac{2\sqrt{3}}{3}$}
\begin{saiten}
\Sp{$\Vec{OH}=s\vec{b}+t\vec{c}$ とおき，垂直条件 $\Vec{AH}\cdot\vec{b}=0,\ \Vec{AH}\cdot\vec{c}=0$ を立てた}{+2}
\Sp{連立方程式を解いて $\Vec{OH}=\dfrac{1}{4}\vec{b}+\dfrac{1}{4}\vec{c}$ を求めた}{+1}
\Sp{$|\Vec{AH}|=\sqrt{3}$ を求め，体積 $V=\dfrac{1}{3}\cdot 2\cdot\sqrt{3}=\dfrac{2\sqrt{3}}{3}$ を求めた}{+2}
\Ssep
\Sm{体積公式の $\dfrac{1}{3}$ を掛け忘れた}{$-1$}
\end{saiten}

\Rh{5}{3}{5}
\Seikai{垂直の証明，$\Pt{MN}=\sqrt{3}$}
\begin{saiten}
\Sp{$\Vec{MN}=\dfrac{1}{2}(\vec{b}+\vec{c}-\vec{a})$ を表した}{+1}
\Sp{$\Vec{MN}\cdot\Vec{OA}=0$ および $\Vec{MN}\cdot\Vec{BC}=0$ を計算して両辺に垂直であることを示した}{+2}
\Sp{$|\Vec{MN}|^{2}=3$ を計算し，$\Pt{MN}=\sqrt{3}$ を求めた}{+2}
\Ssep
\Sm{$\Vec{OA}$ と $\Vec{BC}$ の一方との垂直しか示していない}{$-1$}
\end{saiten}

\Rh{5}{4}{6}
\Seikai{$\Vec{OK}=\dfrac{1}{4}\vec{a}+\dfrac{1}{4}\vec{b}+\dfrac{1}{4}\vec{c}$，$R=\dfrac{\sqrt{7}}{2}$}
\begin{saiten}
\Sp{$\Vec{OK}=p\vec{a}+q\vec{b}+r\vec{c}$ とおき，$|\Vec{OK}|^{2}=|\Vec{OK}-\vec{a}|^{2}$ 等から内積条件 $\vec{a}\cdot\Vec{OK}=2$，$\vec{b}\cdot\Vec{OK}=\vec{c}\cdot\Vec{OK}=\dfrac{5}{2}$ を導いた}{+2}
\Sp{連立方程式を解いて $\Vec{OK}=\dfrac{1}{4}\vec{a}+\dfrac{1}{4}\vec{b}+\dfrac{1}{4}\vec{c}$ を求めた}{+2}
\Sp{$R^{2}=|\Vec{OK}|^{2}=\dfrac{7}{4}$ から外接球の半径 $R=\dfrac{\sqrt{7}}{2}$ を求めた}{+2}
\Ssep
\Sm{外接球の中心 $\Pt{K}$ を平面 $\Pt{OBC}$ 上の点 $\Pt{H}$ と混同した}{$-4$}
\end{saiten}
\betsunote{別解（(3) より直線 $\Pt{MN}$ が辺 $\Pt{OA}$ と辺 $\Pt{BC}$ の共通の垂直二等分線であることを用い，線分 $\Pt{MN}$ の中点 $\Pt{K}$ が $4$ 頂点から等距離にあることを示して求める解法）でも同配点。}

